Let $f:U\to X'$ represent the birational map $X\dashrightarrow X'$ on its largest domain of definition. As in (8.19), properness of $X'$ and the valuative criterion (4.7) show that $U$ contains every codimension-one point of $X$. Hence $\operatorname{codim}(X-U,X)\ge2$.

For any locally free sheaf $\mathcal E$ on $X$, the restriction $\Gamma(X,\mathcal E)\to\Gamma(U,\mathcal E|_U)$ is an isomorphism. Indeed, on an affine open where $\mathcal E$ is free, each coordinate extends uniquely across a subset of codimension at least two by normality and (6.3A).

The differential of $f$ induces $f^*(\bigwedge^q\Omega_{X'/k})\to\bigwedge^q\Omega_{U/k}$ and $f^*(\omega_{X'}^{\otimes n})\to\omega_U^{\otimes n}$. On a dense open subset where $f$ is an isomorphism, both maps are isomorphisms. Since a section of a locally free sheaf on an integral scheme is determined on any dense open subset, the resulting maps on global sections are injective. Using the restriction isomorphisms above gives $h^{q,0}(X')\le h^{q,0}(X)$ and $P_n(X')\le P_n(X)$. Applying the same argument to $X'\dashrightarrow X$ gives the reverse inequalities.
